\documentclass[a4paper,11pt]{article}
\usepackage[utf8]{inputenc}\usepackage[T1]{fontenc}\usepackage[spanish,es-noquoting]{babel}\usepackage{lmodern}
\usepackage[margin=2.2cm]{geometry}\usepackage[hidelinks]{hyperref}\usepackage{booktabs}\usepackage{xcolor}
\setlength{\parindent}{0pt}\setlength{\parskip}{6pt}
\begin{document}
\begin{center}{\LARGE\bfseries Entrega de videos · PFI 2026}\\[4pt]
{\large Análisis de portafolios multiactivo con riesgo geopolítico y recomendaciones explicables mediante IA (Allocia)}\\[4pt]
Huentelaf, Mauricio · Rosso, Máximo · Tutor: Prieto, Leonardo Javier · Curso 551221\end{center}

Los tres videos se adjuntan al correo en formato MP4 y, además, quedan publicados en la plataforma indicada abajo en modo oculto (acceso con el enlace, sin solicitud previa). Se recomienda verlos en el orden numerado.

\begin{tabular}{p{0.6cm}p{6.2cm}p{1.6cm}p{6.6cm}}\toprule
\textbf{N.º} & \textbf{Video} & \textbf{Duración} & \textbf{Enlace}\\\midrule
01 & Problema, investigación y propuesta & 4:54 & \texttt{ENLACE\_01}\\
02 & Arquitectura, implementación y pruebas & 4:59 & \texttt{ENLACE\_02}\\
03 & Allocia explicada: el producto de principio a fin & 4:48 & \texttt{ENLACE\_03}\\\bottomrule
\end{tabular}

\textbf{Contenido.} El video 01 presenta el problema del inversor minorista, la investigación con usuarios (encuesta a 190 personas, tres entrevistas y cinco sesiones de uso), la propuesta, el estado del arte, el modelo de negocio y la viabilidad. El video 02 recorre la arquitectura, el flujo de datos con seis fuentes en vivo, el motor de reglas, la capa de inteligencia artificial y sus controles, la seguridad, la metodología y las pruebas. El video 03 muestra la plataforma tal como la usa un inversor, en la web y en la aplicación móvil.

\textbf{Acceso a la plataforma.} Web: \url{https://allocia.com.ar}. Repositorios: \url{https://github.com/MauriGuth/allocia} (aplicación) y \url{https://github.com/MauriGuth/UADE_PFI} (informe).
\end{document}
